\section{Problem, baseline, and precise contributions}
\label{sec:scope}

The computational objective is to decide whether a planar diagram with
$n$ crossings represents the unknot, with a proved running time of
$n^{O(\log n)}$. The current project has a sound exact recognition
pipeline with several polynomial or restricted-input front ends and
a Bar--Natan scanning implementation of Khovanov homology over $\F$.
The unrestricted Khovanov backend is complete; explicit resource limits
may return \texttt{UNKNOWN}. Its unrestricted worst-case bound is
exponential. This continuation
studies representations that remove two specific kinds of unnecessary
expansion: repeated geometric sheets and repeated algebraic transfer
paths.

\subsection{What was already established}

The audit is anchored to the repository snapshot dated 8 October 2026
at commit \texttt{9cec5275f2d2}, whose full identifier is on the title page
\cite{ProveItSnapshot2026}. The maintained implementation and its tests
are the software baseline. The
following relevant developments precede this report.

\begin{center}
\small
\begin{tabularx}{\textwidth}{@{}p{0.17\textwidth}X@{}}
\toprule
Prior work & Result used or respected here \\
\midrule
Reports 07--09 & Structural certificates, exact component sharing,
specialized braid recognition, certified factorization, and sparse
algebraic filters. \\
Reports 10--12 & Component cobordism contraction, restart accounting,
twist compression, and succinct block cancellation. \\
Report 13 & Exact prescribed homological windows with guard degrees and
support-aware allocation. \\
Reports 14--15 & Continuation quotients, scalar survivor profiles,
nilpotent transfer, and precise geometric assumptions for a sharp
arc-algebra radical bound. \\
Report 16 and integration & Dense and homogeneous coefficient
arithmetic, recovered quantum grading, and its production audit. \\
Unintegrated draft & Complete quantum-triangular survivor transfer,
a binary special contraction, and eight contraction identities. \\
\bottomrule
\end{tabularx}
\end{center}

The last item was found during the continuation audit. Its frozen
source is included as \path{reference/graded_transfer_v1.py}; its
SHA-256 begins \texttt{8c264b87892289c2}. We explicitly treat its
quantum-triangular transfer formula as inherited work
\cite{TriangularDraft2026}. The new algebraic results concern relevant
path support, direction choice, and a sparse scalar representation.
This distinction prevents an independent rediscovery from being
presented as a new project result.

The supplied \texttt{openai/math} catalogue was also inspected
\cite{OpenAIMath2026}. Its general example of releasing manuscripts
together with explicit verification artifacts informs the presentation.
No unverified headline theorem in that catalogue is used as a
mathematical premise here. The arguments below are self-contained
apart from standard covering-space classification, the established
local Khovanov formalism, and the standard homological perturbation
lemma, all cited where used.

\subsection{The two new kernels}

\paragraph{Geometric compression.}
Given a valid dihedral covering presentation of a bordered surface, the
first kernel constructs its complete connected-component topology in
polynomial bit complexity, even if the number of sheets and components
is exponential in the input length. It gives a sharp three-family
classification, rather than returning one record per component.
It also preserves the shared monodromy presentation and answers
individual attachment queries. The result is a concrete special case
of the kind of symbolic topology needed by compressed cutting.

\paragraph{Algebraic compression.}
Given an allocated graded partial complex, the second kernel first
contracts its scalar differential. It then computes the directed graph
of paths that can structurally connect scalar survivors, removes all
other paths, and chooses forward or reverse accumulation for each
remaining component. This preserves every coefficient of the induced
minimal differential. The sparse scalar variant further prevents
unrelated scalar components from being put into a common dense binary
matrix merely because their matching and shifts agree.

These results are new developments relative to the audited project
state. A literature-wide priority claim is not made for standard
covering theory, acyclic path evaluation, or the elementary dihedral
orbit argument. Their relevance is the complete, explicit algorithmic
specialization, its scope, the sharp examples, and the checked code.

\subsection{What the paper does and does not establish}

The mathematical statements distinguish four kinds of assertion:
proved structural theorems, exact algorithmic identities, measured
performance, and prospective applications. In particular, an operation
count for a transfer kernel is not a total running-time bound unless
scalar setup, graph construction, coefficient bit lengths, and output
costs are included. An abstract surface cover is not automatically a
normal surface embedded in a specified knot exterior. The proofs and
interfaces retain these distinctions throughout.

There are two independent routes toward the overall target. The
scanner route needs a bound or a stronger representation for the
surviving complex. The geometric route needs compressed hierarchy
construction and a bound on its restarts. Faster arithmetic in either
route is useful, but it cannot by itself supply the missing structural
bound.

\subsection{Primary topological context}

Bar--Natan's local tangle formalism and cancellation method justify
replacing a partial complex by a homotopy-equivalent complex relative
to its frontier \cite{BarNatan2005,BarNatan2007}. The project's final
Khovanov test relies on the established unknot-detection theorem of
Kronheimer and Mrowka \cite{KronheimerMrowka2011}, with the coefficient
and reduced/unreduced conventions documented in the existing project.
We do not change that test or infer knottedness merely from the size of
a partial complex.

Lackenby's 2026 preprint gives a hierarchy algorithm for
incompressibility and unknot recognition \cite{Lackenby2026}. Section 9
bounds its number of iterations in terms of hierarchy length and
surface pattern complexity. Theorem 14.4 supplies compressed cutting
information, relying on Agol--Hass--Thurston orbit counting
\cite{AHT2006}. The polynomial certificate statements later in the
preprint concern efficiently encoded valid hierarchies; they are not a
quasi-polynomial bound for discovering the needed hierarchy on every
input. Section~\ref{sec:accounting} explains how the present kernels fit
within that distinction.
